\section{Serre functors}
\label{sec:serre}

In this section we review the notion of the Serre functor on $\D(X)$
and then show how to realise the Serre functor on the derived category
$\D(X\times Y)$ using 2-categorical language.

 \subsection{The Serre functor on $\D(X)$.}
\label{subsec:Serretrace}
If $X$ is a space then we consider the functor
\[\SerF\colon\D(X)\to \D(X);\qquad \cE\mapsto \omega_X[\dim
X]\otimes \cE,\] 
where $\omega_X$ is the canonical line bundle of $X$.  Serre duality
then gives natural, bifunctorial isomorphisms
\[ \eta_{\cE,\cF}\colon\Hom_{\D(X)}(\cE, \cF) \stackrel{\sim}{\lra}
\Hom_{\D(X)}(\cF, \SerF \cE)^\chk \]
for any objects $\cE,\cF\in \D(X)$, where ${-}^\chk$ denotes the dual
vector space.

A functor such as $\SerF$, together with isomorphisms as above, was
called a \emph{Serre functor} by Bondal and Kapranov
\cite{BonKapSerre} (see also~\cite{ReiVan}).  From this data, for any
object $\cE\in\D(X)$, define the \emph{Serre trace} as follows:
\[ \Tr\colon\Hom(\cE, \SerF \cE) \ra \k;\qquad \Tr(\alpha) := \eta_{\cE,\cE}(\id_\cE)(\alpha). \]
Note that from this trace we can recover $\eta_{\cE,\cF}$ because
\[ \eta_{\cE,\cF}(\alpha)(\beta)=\Tr(\beta\circ\alpha). \]
We also have the commutativity identity
\[ \Tr(\beta\circ \alpha) = \Tr(\SerF \alpha \circ \beta). \]
Yet another way to encode this data is as a perfect pairing, the
\emph{Serre pairing}:
\[\SP{{-}}{{-}}\colon \Hom(\cE,\cF)\otimes \Hom(\cF,\SerF \cE)\to
\k; \qquad \SP{\alpha}{\beta}:=\Tr(\beta\circ \alpha).\]

 \subsection{Serre kernels and the Serre functor for $\D(X\times
  Y)$.}
\label{subsec:serXY}
We are interested in kernels and the 2-category $\gVar$, so are
interested in Serre functors for product spaces $X\times Y$, and
these have a lovely description in the 2-categorical language.
We can now define one of the key objects in this paper.

\begin{Definition}
For a space $X$, the \emph{Serre kernel} $\SerK_X\in \OHom_\gVar(X,X)$
is defined to be $\Delta_*\omega_X[\dim X]\in\D(X\times X)$, the
kernel inducing the Serre functor on $X$.  Similarly the
\emph{anti-Serre kernel} $\SerK^{-1}_X\in \OHom_\gVar(X,X)$ is defined
to be $\Delta_*\omega^{-1}_X[-\dim X]\in\D(X\times X)$.
\end{Definition}

\begin{Notation}
In string diagrams the Serre kernel will be denoted by a dashed-dotted
line, while the anti-Serre kernel will be denoted by a dashed-dotted
line with a horizontal bar through it (see the pictures in
Section~\ref{subsec:leftrightadjoints}).
\end{Notation}
\noindent
The Serre kernel can now be used to give a natural description, in the
2-category language, of the Serre functor on the product $X\times Y$.

\begin{Proposition}
For spaces $X$ and $Y$ the Serre functor $\SerF_{X\times Y}\colon
\D(X\times Y)\to \D(X\times Y)$
can be taken to be $\SerK_Y\circ{-}\circ\SerK_X$.
\end{Proposition}

\begin{Proof}
The Serre functor on $\D(X\times Y)$ is given by
\[\SerF_{X\times Y}(\Phi)= \Phi\otimes \pi_X^*\omega_X\otimes
\pi_Y^*\omega_Y[\dim X+\dim Y].\]
However, observe that if $\Phi\in \D(X\times Y)$ and $\cE\in \D(X)$
then
\[\Phi\circ \Delta_*\cE\cong \Phi\otimes \pi_X^* \cE,\]
where $\pi_X\colon X\times Y \to X$ is the projection.  This is just a
standard application of the base-change and projection formulas.
Similarly if $\cF\in \D(Y)$ then \( \Delta_*\cF\circ\Phi\cong \pi_Y^*
\cF\otimes\Phi.\) From this the Serre functor can be written as
\[\SerF_{X\times Y}(\Phi)= \SerK_Y\circ\Phi\circ \SerK_X.\qedhere\]
\end{Proof}

\noindent
This means that the Serre trace map on $X\times Y$ is a map
\[ \Tr\colon\THom_\gVar(\Phi, \SerK_Y \circ\Phi\circ\SerK_X) \ra \k \]
which can be pictured as
\[\Tr\left(\pic{Trace}\right) \in \k,\]
where the Serre kernel is denoted by the dashed-dotted line.

We will see below that we have `partial trace' operations which the
Serre trace factors through.


